\documentclass{sig-alternate-05-2015}

\usepackage{tabu}

\let\proof\relax
\let\endproof\relax
\usepackage{amsthm}
\usepackage{caption}
\usepackage{subcaption}
\usepackage{multirow}

\usepackage{algorithm}
\usepackage{algorithmic}

\usepackage{acmcopyright}

\newcommand{\PrimalFig}[2]{
  $2d = #1$
  & \NodeFig{#2}{1}
  & \NodeFig{#2}{2}
  & \NodeFig{#2}{3}
  & \NodeFig{#2}{4}\\
}

\newcommand{\NodeFig}[2]{
   \includegraphics[width=0.085\textwidth]{../../latex/images/run#1node#2.pdf}
}

\usepackage{tkz-graph}
\usepackage{tikz}
\usetikzlibrary{arrows,matrix,decorations.pathreplacing,positioning,chains,fit,shapes,calc} %

\newcommand{\bigoh}{\mathcal{O}}

\DeclareMathOperator{\newdiff}{d} %
\newcommand{\dif}{\newdiff\!}

\newcommand{\vertiii}[1]{{\left\vert\kern-0.25ex\left\vert\kern-0.25ex\left\vert #1
    \right\vert\kern-0.25ex\right\vert\kern-0.25ex\right\vert}}


\DeclareMathOperator*{\argmin}{arg\,min}
\DeclareMathOperator*{\argmax}{arg\,max}
\DeclareMathOperator*{\mini}{minimize}
\DeclareMathOperator*{\maxi}{maximize}
\DeclareMathOperator*{\subto}{\text{subject to}\quad}
\DeclareMathOperator{\BlockDiag}{BlockDiag}
\DeclareMathOperator{\inte}{int}
\newcommand{\1}{\mathbf{1}}
\newcommand{\ld}[1]{{#1}^{[d]}}
\newcommand{\Rh}{\R_{\mathbf{2d}}}
\newcommand{\Sh}{\Sigma_{\mathbf{2d}}}
\newcommand{\Shs}{\Sigma_{\mathbf{2d}}^*}
\newcommand{\din}{d^{-}}
\newcommand{\mdin}{\Delta^{-}}

\newcommand{\la}{\langle}
\newcommand{\ra}{\rangle}

\theoremstyle{definition}
\newtheorem{myprop}{Proposition}
\newtheorem{mytheo}{Theorem}
\newtheorem{mylemm}{Lemma}
\newtheorem{myconj}{Conjecture}
\newtheorem{myprob}{Problem}
\newtheorem{mycons}{Constraint}
\newtheorem{mycor}{Corollary}
\newtheorem{mydef}{Definition}
\newtheorem{myexem}{Example}
\newtheorem{myrem}{Remark}
\newtheorem{myprog}{Program}
\newtheorem{myassu}{Assumption}

\newcommand{\R}{\mathbb{R}}
\newcommand{\Rn}{\R^n}
\newcommand{\Rnn}{\R^{n \times n}}
\newcommand{\Tr}{\top}

\newcommand{\A}{\mathcal{M}}
\newcommand{\bA}{\mathbf{A}}
\newcommand{\G}{G}
\newcommand{\ujsrsos}[1][2d]{\rho_{\mathrm{SOS},#1}(\A)}
\newcommand{\jsrsos}[1][2d]{\rho_{\mathrm{SOS},#1}(\G,\A)}
\newcommand{\jsrsosa}[1][2d]{\rho_{\mathrm{SOS},#1}}
\newcommand{\jsr}{\rho(\A)}
\newcommand{\cjsr}{\rho(\G,\A)}
\newcommand{\cjsrk}[1][k]{\hat{\rho}_{#1}(\G,\A)}
\newcommand{\cjsrp}{\rho(\G',\A')}
\newcommand{\cjsrkp}[1][k]{\hat{\rho}_{#1}(\G',\A')}
\newcommand{\pE}{\widetilde{\mathbb{E}}}
\newcommand{\pmu}{\widetilde{\mu}}

\newcommand{\mapc}{\mathcal{B}}

\newcommand{\Sect}{Section~}
\newcommand{\eqdef}{\triangleq}
\newcommand{\witch}{witched}

\newcommand{\Nodes}{V}
\newcommand{\node}{node}
\newcommand{\nodes}{nodes}
\newcommand{\Arcs}{E}
\newcommand{\arc}{edge}
\newcommand{\arcs}{edges}

\newcommand{\defAG   }{a finite set of matrices $\A$ constrained by an automaton $\G$}
\newcommand{\defAGm  }{a finite set of $m$ matrices $\A$ constrained by an automaton $\G$}
\newcommand{\defAGn  }{a finite set of matrices $\A \subset \mathbb{R}^{n \times n}$ constrained by an automaton $\G$}
\newcommand{\defAGe  }{a finite set of matrices $\A$ constrained by an automaton $\G(\Nodes,\Arcs)$}
\newcommand{\defAGmn }{a finite set of $m$ matrices $\A \subset \mathbb{R}^{n \times n}$ constrained by an automaton $\G$}
\newcommand{\defAGme }{a finite set of $m$ matrices $\A \subset \mathbb{R}$ constrained by an automaton $\G(\Nodes,\Arcs)$}
\newcommand{\defAGne }{a finite set of matrices $\A \subset \mathbb{R}^{n \times n}$ constrained by an automaton $\G(\Nodes,\Arcs)$}
\newcommand{\defAGmne}{a finite set of $m$ matrices $\A \subset \mathbb{R}^{n \times n}$ constrained by an automaton $\G(\Nodes,\Arcs)$}
\newcommand{\defAGie  }{a finite set of invertible matrices $\A$ constrained by an automaton $\G(\Nodes,\Arcs)$}

\DeclareMathOperator{\Diag}{Diag}

\begin{document}

\CopyrightYear{2016} \setcopyright{acmcopyright}
\conferenceinfo{HSCC'16,}{April 12--14, 2016, Vienna, Austria.}
\isbn{978-1-4503-3955-1/16/04}\acmPrice{\$15.00}
\doi{http://dx.doi.org/10.1145/2883817.2883821}


\clubpenalty=10000
\widowpenalty = 10000

\title{Generating Unstable Trajectories for Switched Systems \\
via Dual Sum-Of-Squares Techniques}
\subtitle{[Extended Abstract]}

\numberofauthors{3}
\author{
\alignauthor
Beno\^{i}t Legat\titlenote{This research took place in the framework of the MIT-UCL Internship Program supported by the International Lhoist Berghmans Innovation Chair.}\\
       \affaddr{Universit\'e catholique de Louvain}\\
       \affaddr{4 Av. G. Lema\^itre,}\\
       \affaddr{1348 Louvain-la-Neuve, Belgium}\\
       \email{benoit.legat\\@student.uclouvain.be}
\alignauthor
Rapha\"{e}l M. Jungers\titlenote{R. M. J. is a F.R.S.-FNRS Research Associate. His research is supported by the Belgian Interuniversity Attraction
Poles, and by the ARC grant 13/18-054 from Communaut\'e fran\c{c}aise de Belgique.}\\
       \affaddr{ICTEAM}\\
       \affaddr{Universit\'e catholique de Louvain}\\
       \affaddr{4 Av. G. Lema\^itre,}\\
       \affaddr{1348 Louvain-la-Neuve, Belgium}\\
       \email{raphael.jungers\\@uclouvain.be}
\alignauthor
Pablo A. Parrilo\\
       \affaddr{Laboratory for Information and Decision Systems}\\
       \affaddr{Massachusetts Institute of Technology}\\
       \affaddr{77 Massachusetts Avenue}\\
       \affaddr{Cambridge MA 02139, USA}\\
       \email{parrilo@mit.edu}
}

\maketitle
\begin{abstract}
The joint spectral radius (JSR) of a set of matrices characterizes the
maximal asymptotic growth rate of an infinite product of matrices of
the set.  This quantity appears in a number of applications including
the stability of s\witch{} and hybrid systems.  Many algorithms exist
for estimating the JSR but not much is known about how to generate an
infinite sequence of matrices with an optimal asymptotic growth rate.
To the best of our knowledge, the currently known algorithms select a
small sequence with large spectral radius using brute force (or
branch-and-bound variants) and repeats this sequence infinitely.

In this paper we introduce a new approach to this question, using the
dual solution of a sum of squares optimization program for JSR
approximation. Our algorithm produces an infinite sequence of matrices
with an asymptotic growth rate arbitrarily close to the JSR.  The
algorithm naturally extends to the case where the allowable switching
sequences are determined by a graph or finite automaton.  Unlike the
brute force approach, we provide a guarantee on the closeness of the
asymptotic growth rate to the JSR.  This, in turn, provides new bounds on the
quality of the JSR approximation.
We provide numerical examples illustrating the good performance of the algorithm.
\end{abstract}

\keywords{Joint spectral radius; Sum of squares programming;
S\witch{} Systems; Path-complete Lyapunov functions.}

\section{Introduction}

In recent years, the study of the stability of hybrid systems
has been the subject of extensive research using methods based on classical ideas from Lyapunov theory
and modern mathematical optimization techniques.
Even for s\witch{} linear systems, arguably the simplest class of hybrid systems,
determining stability is undecidable and approximating the maximal asymptotic growth rate is NP-hard \cite{blondel2000boundedness}.
Despite these negative results, the vast range of applications has motivated a wealth of algorithms to approximate
this quantity.

A s\witch{} linear system is characterized by a finite set of matrices $\A \eqdef \{A_1, A_2, \ldots, A_m\} \subset \mathbb{R}^{n \times n}$ and the iteration
\[ x_k = A_{\sigma_k} x_{k-1}, \quad \sigma_k \in \{1, \ldots, m\}. \]
The maximal asymptotic growth rate of this iteration is given by the \emph{joint spectral radius} (JSR).
The JSR $\jsr$ of a finite set of matrices $\A$
is defined as
\[ \jsr = \lim_{k \to \infty} \max_{\sigma \in \{1, \ldots, m\}^k} \|A_{\sigma_k} \cdots A_{\sigma_2} A_{\sigma_1}\|^{1/k}. \]
This definition is independent of the norm used.

The JSR was introduced by Rota and Strang~\cite{rota1960note}
and has many other applications such as wavelets, the capacity of some particular codes, zero-order stability of ordinary differential equations, congestion control in computer networks, curve design and networked and delayed control systems; see \cite{jungers2009joint} for a survey on the JSR and its applications.
Many algorithms exist for estimating the JSR but not much is known on how to generate an infinite
sequence of matrices with an asymptotic growth rate close to the JSR.
However generating such sequence can be of particular interest, depending on the application,
such as exhibiting unstable trajectories for s\witch{} linear systems.
To the best of our knowledge,
the currently known algorithms select a small sequence of matrices with high spectral radius
using brute force (or branch-and-bound variants) and repeats this sequence infinitely \cite{gripenberg1996computing,guglielmi2008algorithm,jungers2014lifted}.

Parrilo and Jadbabaie~\cite{parrilo2008approximation} give an efficient approximation algorithm for computing the JSR using sum of squares (SOS) programming.
In this paper we compute the dual of the corresponding SOS program and
give an algorithm that generates an infinite sequence of matrices with
an asymptotic growth rate arbitrarily close to the JSR.
Unlike the brute force approach,
we provide a guarantee on the quality of the asymptotic growth rate.
As a by-product of the algorithm, the spectral radius of a finite part of
this infinite sequence can be used to give lower bounds on the JSR.
We show on numerical examples that the technique works well in practice.

In some applications the values that $\sigma_k$ can take may depend on $\sigma_{k-1}, \sigma_{k-2}, \ldots$.
These constraints are often conveniently represented using a \emph{finite automaton} and the JSR under such constraints
is called constrained joint spectral radius (CJSR) \cite{dai2012gel}; an example of constrained s\witch{} system is given by Example~\ref{exem:run1} and its automaton is illustrated by Figure~\ref{fig:run}.
Philippe et al.\  generalize the SOS-based approximation algorithm for the CJSR in \cite[\Sect3]{philippe2015stability}.
Our technique is well suited for analysing the more general constrained systems as well.

We provide a new estimate of the accuracy of the SOS-based approximation algorithm for the CJSR which is better
than the previously existing one for sufficiently large SOS degree.
The existing estimate only depends on the dimension of the matrices
while our new one relates the accuracy of the SOS-based approximation algorithm with the combinatorial structure of the automaton representing the constraints.

The automaton representing the constraints can be represented by a
strongly connected labelled directed graph $\G(\Nodes,\Arcs)$,
possibly with parallel \arcs{}.  The labels are elements of the set
$\{1, \ldots, m\}$ and $\Arcs$ is a subset of $\Nodes \times \Nodes
\times \{1, \ldots, m\}$.  We say that $(u,v,\sigma) \in \Arcs$ if
there is an \arc{} between node $u$ and node $v$ with label $\sigma$.

The following will serve as a running example.

\begin{myexem}[Running example\footnote{The source code and instructions are available at the author's web site
to reproduce the numerical results obtained for the running example.
The s\witch{} system considered in Example~\ref{exem:simple1}, Example~\ref{exem:simple2} and Example~\ref{exem:simple3}
is simpler than the running example and the results given in those examples can be obtained by hand.
}]
  \label{exem:run1}
  We borrow the example of \cite[\Sect4]{philippe2015stability}.
  The set of matrices $\A$ is composed of the following four matrices
  \begin{align*}
    A_1 & = A+B
    \begin{pmatrix}
      k_1 & k_2
    \end{pmatrix},&
    A_2 & = A+B
    \begin{pmatrix}
      0 & k_2
    \end{pmatrix},\\
    A_3 & = A+B
    \begin{pmatrix}
      k_1 & 0
    \end{pmatrix},&
    A_4 & = A.
  \end{align*}
  where $k_1 = -0.49$, $k_2 = 0.27$,
  \[
    A =
    \begin{pmatrix}
      0.94 & 0.56\\
      0.14 & 0.46
    \end{pmatrix} \text{ and }
    B =
    \begin{pmatrix}
      0\\
      1
    \end{pmatrix}.
  \]

  The automaton is represented by Figure~\ref{fig:run}.

  \begin{figure}[!ht]
    \centering
    \begin{tikzpicture}[x=40,y=18]
       \Vertex{1}
       \NO[unit=4](1){2}
       \SOEA[unit=2](2){3}
       \EA[unit=2](3){4}
       \tikzset{EdgeStyle/.style = {->}}
       \tikzset{LoopStyle/.style = {->}}
       \tikzset{EdgeStyle/.append style = {bend right=18}}
       \Edge[label=1](1)(3)
       \Edge[label=3](1)(2)
       \Edge[label=2](2)(1)
       \Edge[label=1](2)(3)
       \Edge[label=2](3)(1)
       \Edge[label=3](3)(2)
       \Edge[label=4](3)(4)
       \draw[thick,->] (3) to [out=115,in=65,looseness=10] node [midway, fill=white] {1} (3);
       \Edge[label=1](4)(3)
     \end{tikzpicture}
     \caption{Automaton for the running example.
     The numbers on the \arcs{} are their respective labels.
     }
     \label{fig:run}
  \end{figure}
\end{myexem}


In \Sect\ref{sec:primal}, we explain the SOS-based approximation
algorithm for the CJSR and give our new estimate for its accuracy.
The new bounds explicitly depend on the allowable transitions, through
the graph $G(V,E)$.

In \Sect\ref{sec:dualalgo}, we compute the dual of the corresponding
SOS program and we give our algorithm for generating a sequence of
matrices with an asymptotic growth rate close to the CJSR.

\paragraph{Notations}
We define the automaton $\G^\Tr (\Nodes,\Arcs^\Tr )$ where
\(\Arcs^\Tr  = \{\, (v,u,\sigma) : (u,v,\sigma) \in \Arcs \,\}.\)
We denote as $\Arcs_k$ the subset of $\Arcs^k$ that represents valid paths of length $k$.
The $k$-uple $(\sigma_1, \sigma_2, \ldots, \sigma_k)$ is said to be $\G$-\emph{admissible} if $\sigma_1, \ldots, \sigma_k$ are the respective labels of a path of length $k$.
The \emph{arbitrary switching} case,
that is, when every tuple is $\G$-admissible,
can be seen as the particular case when the automaton has only one node and $m$ self-loops with labels $1, \ldots, m$.
We denote the set of all $k$-uples of $\{1, \ldots, m\}^k$ that are $\G$-admissible as $\G_k$.
The sequence $\sigma_1, \sigma_2, \ldots$ is $\G$-admissible (resp. $\G^\Tr $-admissible) if $(\sigma_1, \ldots, \sigma_k)$ (resp. $(\sigma_k, \ldots, \sigma_1)$) is $\G$-admissible for any $k \geq 1$.
We denote $A_{\sigma_k}\cdots A_{\sigma_1}$ as $A_s$ where $s = (\sigma_1, \ldots, \sigma_k)$ or $s$ is a path with those respective labels.

We denote the set of homogeneous polynomials of degree $2d$ as $\Rh[x]$,
the cone of homogeneous SOS polynomials of degree $2d$ as $\Sh$
and the dual of $\Sh$ as $\Shs$.

\section{Automaton-dependent bounds}
\label{sec:primal}

The definition of the JSR is generalized as follows for constrained systems.
\begin{mydef}[\cite{dai2012gel}]
  \label{def:cjsr}
  The \emph{constrained joint spectral radius} (CJSR) of
  \defAG{}, denoted as $\cjsr$,
  is
  \[ \cjsr = \lim_{k \to \infty}
    \max \Big\{ \|A_s\|^{1/k}\\: s \in \G_k \Big\}.
  \]
\end{mydef}

\subsection{CJSR Approximation via SOS}
\label{sec:sos}
A key result used in the SOS-based approximation algorithm for the unconstrained JSR \cite{parrilo2008approximation} is the following:
\begin{myprop}[{\cite[Proposition~1]{rota1960note}}]
  \label{prop:norm}
  Consider a finite set of matrices $\A$.
  For any $\epsilon > 0$ there exists a norm $\|\cdot\|$ such that
  \[ \|A_\sigma x\| \leq (\jsr+\epsilon)\|x\| \]
  for all $x \in \Rn$ and $\sigma = 1, \ldots, m$.
\end{myprop}

Philippe et al.\  recently gave a generalization of that algorithm, relying on the following extension of Proposition~\ref{prop:norm} to constrained s\witch{} systems.

\vfill\eject
\begin{myprop}[{\cite[Proposition~2.2]{philippe2015stability}}]
  \label{prop:multinorm}
  Consider \defAGe{}.
  For any $\epsilon > 0$, there exists a set of norms $\{\, \|\cdot\|_v : v \in \Nodes \,\}$
  such that
  \[ \|A_\sigma x\|_{v} \leq (\cjsr + \epsilon) \|x\|_{u} \]
  holds for all \arc{} $(u,v,\sigma) \in \Arcs$.
\end{myprop}

Strictly positive homogeneous polynomials are not necessarily norms because their sublevel sets may not be convex.
However we can provide upper bounds for the CJSR using polynomials instead of norms.

\begin{mytheo}
  \label{theo:pos}
  Consider \defAGe{}.
  If there exists a strictly positive homogeneous polynomial $p_v(x)$ of degree $2d$ for all $v \in \Nodes$
  such that
  \[ p_{v}(A_\sigma x) \leq \gamma^{2d}p_{u}(x) \]
  holds for all \arc{} $(u,v,\sigma) \in \Arcs$.
  Then $\cjsr \leq \gamma$.
  \begin{proof}
    Consider a norm $\|\cdot\|$ of $\Rn$ and its corresponding induced matrix norm of $\R^{n \times n}$.
    For each $v \in \Nodes$, we know by compactness of the unit ball in $\Rn$, continuity and strict positivity of $p_v(x)$ that there exist $0 < \alpha_v \leq \beta_v$ such that
    \[ \alpha_v \|x\|^{2d} \leq p_v(x) \leq \beta_v\|x\|^{2d} \]
    for all $x \in \Rn$.
    Let $\alpha = \min_{v \in \Nodes} \alpha_v$ and $\beta = \max_{v \in \Nodes} \beta_v$.

    For a $\G$-admissible $k$-uple $(\sigma_1, \sigma_2, \ldots, \sigma_k)$,
    \[ \|A_{\sigma_k} \cdots A_{\sigma_1}\| = \sup_{x \neq 0} \frac{\|A_{\sigma_k} \cdots A_{\sigma_1}x\|}{\|x\|}. \]
    Consider a path such that the $i$th edge has label $\sigma_i$ for $i = 1, \ldots, k$
    and denote the intermediary nodes of that path as $v_0, v_1, \ldots, v_k$.
    For any $x \in \Rn$, we have
    \begin{align*}
      \|A_{\sigma_k} \cdots A_{\sigma_1}x\| & \leq \alpha_{v_k}^{-\frac{1}{2d}} p_{v_k}(A_{\sigma_k} \cdots A_{\sigma_1} x)^{\frac{1}{2d}}\\
                                            & \leq \alpha_{v_k}^{-\frac{1}{2d}} \gamma p_{v_{k-1}}(A_{\sigma_{k-1}} \cdots A_{\sigma_1} x)^{\frac{1}{2d}}\\
                                            & \leq \alpha_{v_k}^{-\frac{1}{2d}} \gamma^{k} p_{v_0}(x)^{\frac{1}{2d}}
    \end{align*}
    and
    \[ \|x\| \geq \beta_{v_0}^{-\frac{1}{2d}} p_{v_0}(x)^{\frac{1}{2d}} \]
    hence
    \[ \|A_{\sigma_k} \cdots A_{\sigma_1}\| \leq \Big(\frac{\beta_{v_0}}{\alpha_{v_k}}\Big)^{\frac{1}{2d}} \gamma^k \leq \Big(\frac{\beta}{\alpha}\Big)^{\frac{1}{2d}} \gamma^k. \]
    Taking the $k$th root, the limit $k \to \infty$ and using Definition~\ref{def:cjsr} we obtain the result.
  \end{proof}
\end{mytheo}

We relax the positivity condition of Theorem~\ref{theo:pos} by the more tractable sum of squares (SOS) condition and
define $\jsrsos$ as the solution of the following SOS program:
\begin{myprog}[Primal]
  \label{prog:primal}
  \begin{align}
    \notag
    \inf_{p_v(x) \in \Rh[x],\gamma \in \R} \gamma\\
    \label{eq:soscons1}
    p_v(x) & \text{ is SOS} \quad \forall v \in \Nodes\\
    \label{eq:soscons3}
    p_v(x) & \text{ is strictly positive} \quad \forall v \in \Nodes\\
    \notag
    \gamma^{2d} p_u(x) - p_v(A_\sigma x) & \text{ is SOS} \quad \forall (u, v, \sigma) \in \Arcs.
  \end{align}
\end{myprog}

\begin{myrem}
  In practice we can replace \eqref{eq:soscons1} and \eqref{eq:soscons3} by ``$p_v(x) - \epsilon\|x\|_2^{2d}$ is SOS''
  for any $\epsilon > 0$.
  Note that this constrains $p_v(x)$ to be in the interior of the SOS cone,
  which is sufficient for $p_v(x)$ to be strictly positive.
  The bounds given in \Sect\ref{sec:quality} are valid
  if $p_v(x)$ is in the interior of the SOS cone.
\end{myrem}

By Theorem~\ref{theo:pos},
a feasible solution of Program~\ref{prog:primal} gives an upper bound for $\cjsr$,
and thus, for any positive degree $2d$,
\begin{equation}
  \label{eq:upperbound}
  \cjsr \leq \jsrsos.
\end{equation}


\begin{myexem}
  \label{exem:simple1}
  Consider the unconstrained system \cite[Example~2.1]{ahmadi2012joint} with $m=3$:
  \[ \A = \{ A_1 = e_1e_2^\Tr , A_2 = e_2e_3^\Tr , A_3 = e_3e_1^\Tr  \} \]
  where $e_i$ denotes the $i$th canonical basis vector.

  For any $d$, a solution to Program~\ref{prog:primal} is given by
  \[ (p(x), \gamma) = (x_1^{2d} + x_2^{2d} + x_3^{2d}, 1). \]
  Indeed, for example, with $A_1$ we have
  \[ p(e_1e_2^\Tr x) = p(x_2, 0, 0) = x_2^{2d} + 0 + 0 \leq x_1^{2d} + x_2^{2d} + x_3^{2d}. \]
\end{myexem}



\begin{myexem}
  \label{exem:run2}
  Let us reconsider our running example; see Example~\ref{exem:run1}.
  The optimal solution of Program~\ref{prog:primal}
  is represented by
  Figure~\ref{fig:primal} for $2d = 2$, 4, 6, 8, 10 and 12.
  We can see a big difference between the shape of the sublevel sets
  for $2d = 2$ and $2d = 4$
  while between $2d=4$ and $2d=12$,
  the difference seems to be more subtle.
  \begin{figure}[!ht]
    \centering
    \tabulinesep=0.5mm
    \begin{tabu}{{X[0.85,c,m]X[r,m]X[r,m]X[r,m]X[r,m]}}
      & Node 1 & Node 2 & Node 3 & Node 4\\
      \PrimalFig{2}{1}
      \PrimalFig{4}{2}
      \PrimalFig{6}{3}
      \PrimalFig{8}{4}
      \PrimalFig{10}{5}
      \PrimalFig{12}{6}
    \end{tabu}
    \caption{
    Representation of the solutions to Program~\ref{prog:primal}
    with different values of $d$ for the running example.
    The blue curve represents the boundary of the
    1-sublevel set of the optimal
    solution $p_v$ at each node $v \in \Nodes$.
    The dashed curve is the boundary of the unit circle.}
    \label{fig:primal}
  \end{figure}
\end{myexem}

\subsection{Approximation guarantees}
\label{sec:quality}
Philippe et al.\  show, based on \cite{parrilo2008approximation},  the following accuracy guarantees for Program~\ref{prog:primal}.
\begin{mytheo}[{\cite[Theorem~3.6]{philippe2015stability}}]
  \label{theo:matt}
  Consider \defAGn{} and a positive integer $d$.
  The approximation given by Program~\ref{prog:primal} using homogeneous polynomials of degree $2d$
  satisfies:
  \[ \jsrsos \leq {n+d-1 \choose d}^{\frac{1}{2d}} \cjsr. \]
\end{mytheo}

This guarantee only depends on the dimension of the matrices and neither on the number of matrices nor on the automaton.
In this section, we show the following guarantee that relates the accuracy of Program~\ref{prog:primal} to the evolution
of the number of paths in $\G$ for increasing length.
\begin{mytheo}
  \label{theo:second}
  Consider \defAG{} and a positive integer $d$.
  The approximation given by Program~\ref{prog:primal} using homogeneous polynomials of degree $2d$
  satisfies:
  \[ \jsrsos \leq \rho(A(\G))^{\frac{1}{2d}}\cjsr \]
  where $A(\G)$ is the adjacency matrix of $\G$.
\end{mytheo}

\begin{myrem}
  \label{rem:path}
  Consider the matrix norm $\|\cdot\|_\infty$ on $\Rnn$ induced by the infinity norm on $\Rn$.
  It is well known that
  \[ \|A\|_{\infty} = \max_{1 \leq i \leq n} \sum_{j=1}^n |a_{ij}| \]
  where $a_{ij}$ is $(i,j)$ entry of $A$.
  It is also well known that the $(u,v)$ entry of $A(\G)^k$ gives
  the number of paths of length $k$ starting at node $u$ and ending at node $v$ in $\G$.
  Hence $\|A(\G)^k\|_\infty$ is the maximum, over all nodes of $\G$,
  of the number of paths of length $k$ starting at that node.
  By the Gelfand's formula,
  \[
    \rho(A(\G)) = \lim_{k \to \infty} \|A(\G)^k\|_\infty^{1/k}.
  \]
  Therefore the spectral radius $\rho(A(\G))$ represents the ``maximal asymptotic growth rate of the number of paths''.
  This is somewhat related to how constrained the automaton is: the more constrained, the better the approximation.
\end{myrem}

As a corollary of Theorem~\ref{theo:second}, in the trivial cases such that $\rho(A(\G)) = 1$, the approximation is exact.
This corresponds to the case where every node of $\G$ has indegree and outdegree 1.
In that case, the graph forms a cycle of some length $k$ and the CJSR is simply the $k$th root of the spectral radius of the product of the matrices along this cycle.

The results of Theorem~\ref{theo:matt}, Theorem~\ref{theo:second} and \eqref{eq:upperbound} are summarized by the following corollary.
\begin{mycor}
  \label{cor:mix}
  Consider \defAGn{} and a positive integer $d$,
  the approximation given by Program~\ref{prog:primal} using homogeneous polynomials of degree $2d$
  satisfies:
  \begin{multline*}
    \min\Big\{{n+d-1 \choose d}, \rho(A(\G))\Big\}^{-\frac{1}{2d}} \jsrsos\\\leq \cjsr \leq \jsrsos.
  \end{multline*}
  where $A(G)$ is the adjacency matrix of the automaton $\G$.
\end{mycor}
We see that we can have arbitrary accuracy by increasing $d$.

For the arbitrary switching case, $\rho(A(\G)) = m$.
Theorem~\ref{theo:second} was already known in this particular case \cite[Theorem~4.3]{parrilo2008approximation}.
This section can thus be seen as a generalization of \cite[\Sect4]{parrilo2008approximation}.

Our proof technique relies on the analysis of an iteration in the vector space of polynomials of degree $2d$.
When this iteration converges, it converges to a feasible solution of Problem~\ref{prog:primal}.
By analysing this iteration as affine iterations in this vector space, we derive a sufficient condition for its convergence
and thus an upper bound for $\jsrsos$.

Consider the iteration
\begin{align}
  \notag
  P_{v,0}(x) & = 0,\\
  \label{eq:iter}
  P_{v,k+1}(x) & = Q_v(x) + \frac{1}{\beta}
  \sum_{(u,v,\sigma) \in \Arcs} P_{u,k}(A_\sigma x), \quad  v \in \Nodes
\end{align}
for fixed homogeneous polynomial $Q_v(x)$ of degree $2d$ in $n$ variables (not necessarily different)
and a constant $\beta > 0$.

When this iteration converges, it converges to a feasible solution of Program~\ref{prog:primal}.
\begin{mytheo}
  \label{theo:conv}
  Consider a constant $\beta > 0$.
  If there exists homogeneous polynomials $Q_v(x)$ in the interior of the SOS cone such that
  the iteration \eqref{eq:iter} converges then
  \[
    \jsrsos \leq \beta^{\frac{1}{2d}}.
  \]
  \begin{proof}
    Suppose the iteration converges to the polynomials $P_{v,\infty}(x)$.
    It is easy to show by induction that $P_{v,k}(x)$ is SOS for all $k$.
    It is trivial for $k = 0$ and if it is true for $k$ then it is also true for $k+1$ by \eqref{eq:iter}.
    Since the SOS cone is closed, $P_{v,\infty}$ is SOS.
    Now by \eqref{eq:iter}, for each $v \in \Nodes$,
    \[ P_{v,\infty}(x) = Q_v(x) + \frac{1}{\beta}
    \sum_{(u,v,\sigma) \in \Arcs} P_{u,\infty}(A_\sigma x) \]
    so $P_{v,\infty}(x)$ is also in the interior of the SOS cone.
    For each edge $(u,v,\sigma)$, by manipulating the above equation, we have
    \[ \beta P_{v,\infty}(x) - P_{u,\infty}(A_\sigma x) = \beta Q_v(x) +
    \sum_{\substack{(u',v,\sigma') \in \Arcs\\(u',\sigma') \neq (u, \sigma)}} P_{u',\infty}(A_{\sigma'} x) \]
    so $\beta P_{v,\infty}(x) - P_{u,\infty}(A_\sigma x)$ is in the interior of the SOS cone.
    Therefore $(\{\, P_{v,\infty}(x) : v \in \Nodes \,\}, \beta^{\frac{1}{2d}})$ is a feasible solution of Program~\ref{prog:primal}.
  \end{proof}
\end{mytheo}

\begin{myrem}
  Before going into the details, we first provide an intuition on how we
  prove Theorem~\ref{theo:second} using the iteration~\eqref{eq:iter}.
  Consider the following iteration (it is not the same as
  \eqref{eq:iter} but the sufficient condition of convergence that we
  develop in Theorem~\ref{theo:cond} is the same for both iterations)
  \begin{align*}
    P_{v,0}(x) & = Q_v(x),\\
    P_{v,k+1}(x) & = \frac{1}{\beta}
    \sum_{(u,v,\sigma) \in \Arcs} P_{u,k}(A_\sigma x), \quad  v \in \Nodes.
  \end{align*}
  We can see that
  \begin{equation}
    P_{v,k+1}(x) = \frac{\cjsr^{2dk}}{\beta^k}\sum_{\substack{s \in \Arcs_k,\\s(k+1)=v}} Q_{s(1)}\Big(\frac{A_s}{\cjsr^{k}} x\Big)
  \end{equation}
  where $s(i)$ denotes the $i$th \node{} of the path $s$.
  As $k \to \infty$, $A_s / \cjsr^{k}$ tends to zero.
  However, the number of terms in the right-hand side may tend to infinity.
  It is equal to the number of paths of length $k$ whose last node is $v$.
  Therefore $\beta/\cjsr^{2d}$ should be at least the asymptotic value of the $k$th root of the number of paths of length $k$.
  We have seen in Remark~\ref{rem:path} that this value is $\rho(A(\G))$.
\end{myrem}

In view of Theorem~\ref{theo:conv}, it is thus natural to analyse
under which condition the iteration~\ref{eq:iter} converges. While we
could convert the idea behind Remark 3 into a full proof, we present
here a slightly different approach. Recall that
iteration~\ref{eq:iter} is an affine map on the vector space of
homogeneous polynomials of degree $2d$.  We would like to choose a
basis for this vector space to represent the linear operator as a
matrix and use a sufficient condition for convergence depending on the
spectral radius of the matrix.  A simple choice is the monomial basis
$x^\alpha$ with $\alpha$ such that $|\alpha| = d$ where $x^\alpha$
denotes $x_1^{\alpha_1}\cdots x_n^{\alpha_n}$ and $|\alpha|$ denotes
$\alpha_1 + \cdots + \alpha_n$.  We choose instead a basis for which
each monomial $x^\alpha$ is scaled by the coefficient ${d \choose
  \alpha} = \frac{d!}{\alpha_1!\alpha_2!\cdots\alpha_n!}$ so that the
corresponding lift operation for vectors of $\R^n$ that we define
below preserves the Euclidean norm (up to power).

\begin{mydef}[$d$-lift]
  Let $N_d = {n+d-1 \choose d}$.
  The $d$-lift of a vector $x \in \R^n$ is the vector $x^{[d]} \in \R^{N_d}$
  (indexed by the $n$-uples $\alpha \in \mathbb{N}^n$
  such that $|\alpha| = d$)
  such that
  \( x_\alpha^{[d]} = {d \choose \alpha} x^\alpha. \)
\end{mydef}

For any matrix $A \in \R^{n \times n}$, the $d$-lift induces an associated map $A^{[d]} \in \R^{N_d \times N_d}$
which is the unique $N_d \times N_d$ matrix that satisfies $(Ax)^{[d]} = A^{[d]}x^{[d]}$.
We also denote $\A^{[d]} \eqdef \{A_1^{[d]}, \ldots, A_m^{[d]}\}$.

As mentioned above, we have chosen to scale the monomials so that $\|x^{[d]}\|_2 = \|x\|_2^d$ where $\|\cdot\|_{2}$ is the Euclidean norm.
The purpose of this choice is to have the following lemma for the lift operation.
\begin{mylemm}[{\cite[Lemma~3.5]{philippe2015stability}}]
  \label{lemm:Ad}
  Consider \defAG{} and a positive integer $d$.
  The following identity holds:
  \[ \rho(\G, \A^{[d]}) = \cjsr^d. \]
\end{mylemm}

We can now formulate a sufficient condition for the convergence of the iteration~\ref{eq:iter}.
\begin{mytheo}
  \label{theo:cond}
  Let
  \[ B_{2d} = \sum_{(u,v,\sigma) \in \Arcs} (e_ue_v^\Tr ) \otimes A_\sigma^{[2d]} \]
  where $\otimes$ is the Kronecker product.
  If
  \( \beta < \rho(B_{2d}) \)
  then the iteration~\eqref{eq:iter} converges
  for any homogeneous polynomials $Q_v(x)$ of degree $2d$ in the interior of the SOS cone.
  \begin{proof}
    We can rewrite the polynomials in terms of the basis of scaled monomials.
    For each $v \in \Nodes$, we consider the vector $q_v \in \R^{N_{2d}}$ such that $Q_v(x) = \langle q_v, x^{[2d]} \rangle$.
    We see that $Q_v(A_\sigma x) = \langle q_v, A^{[2d]}_\sigma x^{[2d]} \rangle = \langle (A^{[2d]}_\sigma)^\Tr q_v,  x^{[2d]} \rangle$.
    The iteration~\eqref{eq:iter} is
    \begin{align}
      \notag
      p_{v,0} & = 0,\\
      \notag
      p_{v,k+1} & = q_j + \frac{1}{\beta} \sum_{u \in \Nodes} \Big(\sum_{(u,v,\sigma) \in \Arcs} A_\sigma^{[2d]}\Big)^\Tr p_{u,k}, \quad v \in \Nodes
    \end{align}
    which can be rewritten as
    \begin{equation}
      \label{eq:compactit}
      p_{0} = 0, \quad p_{k+1} = q + \frac{1}{\beta} B_{2d}^\Tr p_k.
    \end{equation}
    where
    the vectors $p_k \in \R^{|\Nodes|N_{2d}}$ are (resp. $q \in \R^{|\Nodes|N_{2d}}$ is) the concatenation of $p_{1,k}, \ldots, p_{|\Nodes|,k}$ (resp. $q_1, \ldots, q_{|\Nodes|}$).


    The iteration in Equation \eqref{eq:compactit} is an affine iteration, and it is well known that it converges if $\rho(B_{2d}) < \beta$,
    independently of $q$.
  \end{proof}
\end{mytheo}

\begin{myrem}
  If $\rho(B_{2d}) < \beta$, $p_\infty \eqdef \lim_{k\to\infty} p_k$ can be computed by solving the linear system $(I - B_{2d}/\beta)^\Tr p_\infty = q$.
\end{myrem}

As a corollary of Theorem~\ref{theo:conv} and Theorem~\ref{theo:cond},
\begin{equation}
  \label{eq:convcond}
  \jsrsos \leq \rho(B_{2d})^{\frac{1}{2d}}.
\end{equation}
To prove Theorem~\ref{theo:second}, it remains to find a relation between $\rho(B_{2d})$ and $\cjsr$.

Given a sub-multiplicative norm $\|\cdot\|$ on $\R^{n_1 \times n_1}$,
consider the sub-multiplicative norm $\vertiii{\cdot}$ on $\R^{n_2n_1 \times n_2n_1}$ induced by $\|\cdot\|$ defined in \cite{kozyakin2014berger} as
\begin{equation}
  \label{eq:blocknorm}
  \vertiii{M} \eqdef \max_{1 \leq i \leq N} \sum_{j=1}^{n_2} \|m_{ij}\|
\end{equation}
for a block $n_2 \times n_2$ matrix $M = (m_{ij})_{i,j=1}^{n_2}$ with $m_{ij} \in \R^{n_1 \times n_1}$.
We can show that
\begin{mylemm}
  \label{lemm:cs}
  Given a norm $\|\cdot\|$ on $\R^{n_1 \times n_1}$ and the norm $\vertiii{\cdot}$ defined by \eqref{eq:blocknorm},
  for any matrices $A_1, \ldots, A_m \in \R^{n_1 \times n_1}$ and any nonnegative matrices $B_1, \ldots, B_m \in \R^{n_2 \times n_2}$,
  we have
  \[
    \vertiii{\sum_{k=1}^m B_k \otimes A_k}
    \leq \bigg\|\sum_{k=1}^m B_k\bigg\|_\infty \max_{1 \leq k \leq m} \|A_k\|.
  \]
  \begin{proof}
    Using \eqref{eq:blocknorm} and the subadditivity of the norm $\|\cdot\|$,
    \begin{align*}
      \vertiii{\sum_{k=1}^m B_k \otimes A_k}
      & = \max_{1 \leq i \leq N} \sum_{j=1}^N \left\|\sum_{k=1}^m (B_k)_{ij} A_k\right\|\\
      & \leq \max_{1 \leq i \leq N} \sum_{j=1}^N \sum_{k=1}^m (B_k)_{ij} \|A_k\|\\
      & \leq \big(\max_{1 \leq k \leq m} \|A_k\|\big) \max_{1 \leq i \leq N} \sum_{j=1}^N \sum_{k=1}^m (B_k)_{ij}\\
      & \leq \left\|\sum_{k=1}^m B_k\right\|_\infty \max_{1 \leq k \leq m} \|A_k\|.
    \end{align*}
  \end{proof}
\end{mylemm}

This Lemma allows us to prove the following generalization of~\cite[Lemma~1]{blondel2005computationally}.
\begin{mylemm}
  \label{lemm:srsjsr}
  For any finite set of matrices $\A = \{A_1, \ldots, A_m\}$ of $\R^{n_1 \times n_1}$, any nonnegative matrices $B_1, \ldots, B_m \in \R^{n_2 \times n_2}$ and any automaton $\G(\Nodes,\Arcs)$ such that for each path $p$ that is not $\G$-admissible, $B_p = 0$,
  we have
  \[
    \rho(\sum_{i=1}^m B_i \otimes A_i) \leq \rho(\sum_{i=1}^m B_i) \rho(\G, \A).
  \]
  \begin{proof}
    By Lemma~\ref{lemm:cs},
    \begin{align*}
      \vertiii{\Big(\sum_{i=1}^m B_i \otimes A_i\Big)^k}
      & = \vertiii{\sum_{s \in \Arcs^k} B_p \otimes A_s}\\
      & = \vertiii{\sum_{s \in \Arcs_k} B_p \otimes A_s}\\
      & \leq \Big\|\sum_{s \in \Arcs_k} B_p\Big\|_\infty \max_{s \in \Arcs_k} \|A_s\|\\
      & = \Big\|\Big(\sum_{i=1}^m B_i\Big)^k\Big\|_\infty \max_{s \in \Arcs_k} \|A_s\|.
    \end{align*}
    Taking the $k$th root and the limit of $k \to \infty$ we obtain the result.
  \end{proof}
\end{mylemm}

We are now ready to prove Theorem~\ref{theo:second}.

\begin{proof}[Proof of Theorem~\ref{theo:second}]
  We have
  \begin{align*}
    \jsrsos
    & \leq \rho(B_{2d})^{\frac{1}{2d}}\\
    & \leq \Bigg[\rho\Big(\sum_{(i,j,\sigma) \in \Arcs} e_ie_j^\Tr \Big) \rho(\G,\A^{[2d]})\Bigg]^{\frac{1}{2d}}\\
    & = [\rho(A(\G)) \rho(\G,\A^{[2d]})]^{\frac{1}{2d}}\\
    & = \rho(A(\G))^{\frac{1}{2d}} \rho(\G,\A).
  \end{align*}
  where the first inequality is \eqref{eq:convcond},
  the second is given by Lemma~\ref{lemm:srsjsr},
  the first equality is given by the fact that
  \[ \sum_{(i,j,\sigma) \in \Arcs} e_ie_j^\Tr  = A(\G) \]
  and the second equality is given by Lemma~\ref{lemm:Ad}.
\end{proof}

\section{Finding high-growth sequences}
\label{sec:dualalgo}
\subsection{Dual SOS program}
\label{sec:dual}
Consider the feasibility version of Program~\ref{prog:primal},
i.e. $\gamma$ is fixed and we want to find polynomials such that the
constraints hold.

The dual variables are in the dual of the SOS cones,
i.e. they are linear functionals over homogeneous polynomials of degree $2d$
that gives a positive value when applied to a SOS polynomials.
Barak et al.~\cite{barak2012hypercontractivity} introduced the \emph{pseudo-expectation} $\pE$ as a way to interpret those linear functionals.
It can be seen as the expectation of a \emph{pseudo-distribution} $\pmu$.
That is,
\begin{equation}
  \label{eq:exp}
  \la \pE, p(x) \ra = \pE[p(x)] = \int_{\mathbb{S}^{n-1}} p(x) \dif \pmu.
\end{equation}
We say that $\pmu$
is only a `pseudo-distribution' (and not a classical distribution) in the sense that it
is not necessarily a positive measure,
it might be possible that $\pE[p(x)] < 0$ while $p(x)$ is nonnegative.
However, if $p(x)$ is an homogeneous SOS polynomial of degree $2d$
then $\pE[p(x)] \geq 0$ since $\pE$ belongs to the dual of the SOS cone.

Note that for any matrix $A \in \R^{n \times n}$, the polynomial $p(A x)$ depends linearly on $p(x)$.
Let $\mathcal{A}$ be the linear map such that $\mathcal{A}(p(x)) = p(Ax)$ for all $p(x) \in \R[x]_{\mathbf{2d}}$.
We define the adjoint map $\mathcal{A}^*$,
the linear map such that $\la \mathcal{A}^* \pE, p(x) \ra = \la \pE, \mathcal{A}(p(x)) \ra$ for all $\pE \in (\R[x]_{\mathbf{2d}})^*$ and $p(x) \in \R[x]_{\mathbf{2d}}$.

The dual of the feasibility version of the primal is given by
\begin{myprog}[Dual]
  \label{prog:dual}
  \begin{align}
    \label{eq:dualconsexp1}
    \sum_{(u,v,\sigma) \in E} \mathcal{A}_\sigma^* \pE_{uv\sigma} & \succeq \gamma^{2d} \sum_{(v,w,\sigma) \in \Arcs} \pE_{vw\sigma}, \forall v \in \Nodes,\\
    \notag
    \pE_{uv\sigma} & \in \Shs, \qquad \forall (u,v,\sigma) \in \Arcs,\\
    \label{eq:dualconsexp3}
    \sum_{(u,v,\sigma) \in E} \pE_{uv\sigma}\Big[\sum_{i=1}^n x_i^{2d}\Big] & = 1.
  \end{align}
\end{myprog}

Note that the dual constraint \eqref{eq:dualconsexp1} is equivalent to:
\begin{equation}
  \label{eq:dualconsexp}
  \sum_{(u,v,\sigma) \in \Arcs} \pE_{uv\sigma}[p(A_\sigma x)] \geq \gamma^{2d} \sum_{(v,w,\sigma) \in \Arcs} \pE_{vw\sigma}[p(x)]
\end{equation}
for all $v \in \Nodes$ and for all $p(x) \in \Sigma_{\mathbf{2d}}$.

\begin{myexem}
  \label{exem:simple2}
  Consider Example~\ref{exem:simple1}.
  For $i = 1, 2, 3$, let $\pE_i$ be the dual solution at the \arc{} corresponding to the matrix $A_i$.
  For any $d$, we can see that the dual solution for $\gamma = 1$ is such that
  the only monomial $x^\alpha$ such that $\pE_1[x^\alpha]$ (resp. $\pE_2[x^\alpha]$, $\pE_3[x^\alpha]$) is non-zero is $x_1^{2d}$ (resp. $x_2^{2d}$, $x_3^{2d}$)
  and $\pE_1[x_1^{2d}] = \pE_2[x_2^{2d}] = \pE_3[x_3^{2d}] = 1/3$.
  Note that it means that $\pmu_1 = \delta_{(1,0,0)}/3$, $\pmu_2 = \delta_{(0,1,0)}/3$ and $\pmu_3 = \delta_{(0,0,1)}/3$
  where $\delta_x$ is the Dirac measure centered on $x$.
  Since there also exists a primal solution for $\gamma = 1$ (see Example~\ref{exem:simple1}) we have $\jsrsos = 1$ for any $d$.
\end{myexem}

\begin{myexem}
  \label{exem:run3}
  We continue the running example; see Example~\ref{exem:run1} and Example~\ref{exem:run2}.

  For all $d$, $\pE_{212} = \pE_{323} = \pE_{344} = \pE_{431} = 0$
  hence the node 4 is ``unused'' by the dual.
  For $2d = 2, 4, 6, 8$, $\pE_{123} = \pE_{231} = 0$ so
  the node $2$ is ``unused'' for low degree.

  At first, one could think that the dual variables can be used to reduce the systems,
  e.g. remove nodes or edges.
  However, as we will see, it would be a mistake to remove the node 2.

  For $2d = 10$, $\pE_{123}$ and $\pE_{231}$ are not zero and are of the ``same order or magnitude'' than $\pE_{131}$ and $\pE_{312}$.
  Then for $2d = 12$, $\pE_{123}$ and $\pE_{231}$ have ``larger magnitude'' than $\pE_{131}$ and $\pE_{312}$.
  This observation will be useful for Example~\ref{exem:run4}.

  We can see that while the shape of the primal variables change a lot between $2d = 2$ and $2d = 4$ as mentioned in Example~\ref{exem:run2},
  the ``important'' change for the dual variables happens around $2d = 10$.

  It is also interesting to notice that the matrices corresponding to the dual variables have low rank.
  For example, for $2d = 2$, the measure corresponding to the dual variable $\pE_{131}$
  (resp. $\pE_{312}$, $\pE_{331}$) is the Dirac measure $0.324 \times \delta_{(0.917,0.399)}$
  (resp. $0.229 \times \delta_{(0.875,0.485)}$, $0.447 \times \delta_{(0.757,-0.653)}$).
\end{myexem}

\subsection{Basic algorithm}
\label{sec:seq}
In this section we give an algorithm that generates an infinite sequence of matrices such that
the asymptotic value of the norm of the product of the matrices is arbitrarily close to the CJSR.
Note that by Definition~\ref{def:cjsr}, this asymptotic value must be smaller than the CJSR.

\begin{myrem}
  The algorithm generates the sequence of matrices in the \emph{reverse order}.
  We can either see the sequence as $\G$-admissible and left infinite
  or $\G^\Tr $-admissible and right infinite.
\end{myrem}

Denote the indegree of a node $v \in \Nodes$ as $\din(v)$ and the
maximum indegree of $\G$ as $\mdin(\G) = \max_{v \in \Nodes} \din(v)$.
Suppose that we are given a polynomial $p_0(x) \in \inte(\Sh)$,
i.e. in the interior of the cone of SOS homogeneous polynomials of degree $2d$.
Since it is in the interior, $\pE_{uv\sigma}[p_0(x)] > 0$ for all $(u,v,\sigma) \in \Arcs$ such that $\pE_{uv\sigma}$ is not zero.
By \eqref{eq:dualconsexp3}, at least one dual variable is nonzero and we can choose one of such \arcs{}, say $(v_1,v_0,\sigma_0)$.

The design of the algorithm is to
build a $\G^\Tr $-admissible sequence $(v_1,v_0,\sigma_0), (v_2,v_1,\sigma_1), \ldots$
such that
\[
  \theta_k \eqdef \pE_{v_{k+1}v_k\sigma_k}[p_{0}(A_{\sigma_1} \cdots A_{\sigma_k}x)]
\]
remains large for increasing $k$.
Indeed, this would suggest that the product $A_{\sigma_1} \cdots A_{\sigma_k}$ has a large norm.
The key idea is that maintaining a large value for $\theta_k$ is possible thanks to \eqref{eq:dualconsexp}.

  \begin{algorithm}
    \caption{Producing a sequence of matrices with an asymptotic growth rate close to the CJSR.}
    \label{algo:prod}
    \begin{algorithmic}
      \STATE{Pick an arbitrary polynomial $p_0(x) \in \inte(\Sh)$}
      \STATE{Pick an \arc{} $(v_1,v_0,\sigma_0) \in \Arcs$ such that $\pE_{v_1v_0\sigma_0}[p_0(x)] > 0$}
      \STATE{Set $p_1(x) \leftarrow p_0(x)$}
      \FOR{$k = 1,2,\ldots$}
      \STATE $(v_{k+1},v_k,\sigma_k) \leftarrow \argmax_{(u,v_k,\sigma) \in \Arcs} \pE_{uv_k\sigma}[p_k(A_\sigma x)]$
      \STATE $p_{k+1} \leftarrow p_k(A_{\sigma_k} x)$
      \ENDFOR
    \end{algorithmic}
  \end{algorithm}

Expanding $p_{k+1}(x)$, we obtain
\[ p_{k+1}(x) = p_0(A_{\sigma_1}A_{\sigma_2} \cdots A_{\sigma_k} x). \]

The dual constraint \eqref{eq:dualconsexp} enables us to guarantee the following.

\begin{mylemm}
  \label{lemm:maxdegpath}
  Consider \defAGe{}.
  For any positive integer $d$,
  using Program~\ref{prog:dual} with any $\gamma < \jsrsos$,
  Algorithm~\ref{algo:prod} produces a $\G^\Tr $-admissible sequence
  $(v_1,v_0,\sigma_0), (v_2,v_1,\sigma_1), \ldots$ %
  that satisfies the following inequality for all $k \geq 1$:
  \begin{equation*}
    \theta_{k}
    \geq \frac{\gamma^{2d}}{\din(v_k)}
    \theta_{k-1}.
  \end{equation*}
  \begin{proof}
    For any integer $k \geq 1$, since $p_0(x)$ is SOS,
    \[ p_0(A_{\sigma_1} \cdots A_{\sigma_{k-1}}x) \text{ is SOS.} \]
    Therefore, by \eqref{eq:dualconsexp},
    \begin{multline*}
      \sum_{(u,v_k,\sigma) \in \Arcs} \pE_{uv_k\sigma}[p_0(A_{\sigma_1} \cdots A_{\sigma_{k-1}} A_\sigma x)]\\
      \geq \gamma^{2d} \sum_{(v_k,w,\sigma) \in \Arcs} \pE_{v_kw\sigma}[p_0(A_{\sigma_1} \cdots A_{\sigma_{k-1}} x)]
    \end{multline*}
    Since the dual variables $\pE_{v_kw\sigma}$ of the right-hand side are in the dual of the SOS cone,
    all the terms of the right-hand side are positive.
    Since one of them is $\theta_{k-1}$,
    \[
      \sum_{(u,v_k,\sigma) \in \Arcs} \pE_{uv_k\sigma}[p_0(A_{\sigma_1} \cdots A_{\sigma_{k-1}} A_\sigma x)]\\
      \geq \gamma^{2d} \theta_{k-1}.
    \]
    Since the left-hand side has $\din(v_k)$ positive terms and
    Algorithm~\ref{algo:prod} picks the term with highest value, the result follows.
  \end{proof}
\end{mylemm}

To give a guarantee for the norm of the product of the matrices of the sequence,
we introduce the following two Lemmas.

\begin{mylemm}
  \label{lemm:svd}
  For any matrix $A \in \R^{n \times n}$ and symmetric positive semidefinite matrix $Q$,
  the following inequality holds
  \[ \rho(A^\Tr QA) \leq \rho(Q) \rho(A^\Tr A). \]
\end{mylemm}

\begin{mylemm}
  \label{lemm:sosbound}
  For any polynomial $p(x) \in \inte(\Sh)$ and any matrix $A \in \Rnn$, there exists a constant $\beta$ that does not depend on $A$
  such that

  \[ \beta\|A\|_2^{2d} p(x) - p(Ax) \quad \text{is SOS}. \]
  Moreover we can choose $\beta = \kappa(Q)$ where $Q \in \mathcal{S}^{{n+d \choose d}}$
  is the symmetric matrix such that $p(x) = (x^{[d]})^\Tr Qx^{[d]}$
  and $\kappa(Q) = \rho(Q) \rho(Q^{-1})$ is the condition number of $Q$.
  \begin{proof}
    Consider the matrix $Q$ defined in the statement of the lemma.
    Note that since $p(x) \in \inte(\Sh)$, $Q$ is positive definite.
    We can see that
    \[ ((Ax)^{[d]})^\Tr Q(Ax)^{[d]} = (\ld{x})^\Tr (A^{[d]})^\Tr Q\ld{A}\ld{x}. \]
    Moreover, for all $x \in \Rn$,
    \[ (\ld{x})^\Tr Q\ld{x} \geq \|\ld{x}\|_2/\rho(Q^{-1}) \]
    and by Lemma~\ref{lemm:svd},
    \[ (\ld{x})^\Tr (A^{[d]})^\Tr Q\ld{A}\ld{x} \leq \rho(Q)\rho((\ld{A})^\Tr \ld{A})\|\ld{x}\|_2. \]
    Using the fact that $\rho(A^{[d]}) = \rho(A)^d$ and $(AB)^{[d]} = A^{[d]}B^{[d]}$ for any matrix $A$ and $B$,
    we have
    \[ \kappa(Q) \rho(A^\Tr A)^{d} Q - (A^{[d]})^\Tr QA^{[d]} \succeq 0. \]
  \end{proof}
\end{mylemm}

\begin{mytheo}
  \label{theo:third}
  Consider \defAGie{}.
  For any positive integer $d$,
  using Program~\ref{prog:dual} with any $\gamma < \jsrsos$,
  Algorithm~\ref{algo:prod} produces a $\G^\Tr $-admissible sequence
  $(v_1,v_0,\sigma_0),\allowbreak(v_2,v_1,\sigma_1), \ldots$ %
  such that
  \begin{equation}
    \lim_{k \to \infty} \|A_{\sigma_1} \cdots A_{\sigma_k}\|_2^{\frac{1}{k}} \geq \frac{\gamma}{[\mdin]^{\frac{1}{2d}}}.
  \end{equation}
  \begin{proof}
    Consider the sequence produced by Algorithm~\ref{algo:prod}.
    Since the set of \arcs{} $\Arcs$ is finite,
    there must be an \arc{} $(\bar{u},\bar{v},\bar{\sigma}) \in \Arcs$ that appears infinitely many times in the sequence.
    Let $k_1$ be the smallest integer such that $(v_{k_1+1},v_{k_1},\sigma_{k_1}) = (\bar{u},\bar{v},\bar{\sigma})$ and
    let $q(x) = p_{k_1+1}(x)$.
    Since the matrices of $\A$ are invertible and $p_0(x) \in \inte(\Sh)$, we know that $q(x) \in \inte(\Sh)$.

    For an arbitrarily large integer $K$, there exists an integer $k \geq K$ such that $(v_{k_1+k+1},v_{k_1+k}, \sigma_{k_1+k}) = (\bar{u},\bar{v},\bar{\sigma})$.
    Let $s_K = (\sigma_{k_1+k}, \ldots, \sigma_{k_1+1})$.
    By Lemma~\ref{lemm:maxdegpath}, we have
    \begin{equation}
      \label{eq:cycling}
      \pE_{\bar{u}\bar{v}\bar{\sigma}}[q(A_{s_K} x)] \geq \frac{\gamma^{2dk}}{[\mdin]^{k}} \pE_{\bar{u}\bar{v}\bar{\sigma}}[q(x)].
    \end{equation}
    By Lemma~\ref{lemm:sosbound}, there exists a constant $\beta$ that does not depend on $A_{s_K}$ such that%
    \[ \beta \|A_{s_K}\|_2^{2d}q(x) - q(A_{s_K} x) \text{ is SOS}. \]
    Therefore,
    \[ \pE_{\bar{u}\bar{v}\bar{\sigma}}[\beta \|A_{s_K}\|_2^{2d}q(x)] \geq \pE_{\bar{u}\bar{v}\bar{\sigma}}[q(A_{s_K} x)] \]
    and by \eqref{eq:cycling},
    \[ \beta \|A_{s_K}\|_2^{2d}\pE_{\bar{u}\bar{v}\bar{\sigma}}[q(x)] \geq \frac{\gamma^{2dk}}{[\mdin]^k} \pE_{\bar{u}\bar{v}\bar{\sigma}}[q(x)]. \]
    Since $\pE_{\bar{u}\bar{v}\bar{\sigma}}[q(x)] > 0$,
    \[ \frac{\gamma^{2dk}}{[\mdin]^k} \leq \beta \|A_{s_K}\|_2^{2d} \]
    or equivalently
    \[ \frac{\gamma}{[\beta^{\frac{1}{k}} \mdin]^{\frac{1}{2d}}} \leq \|A_{s_K}\|_2^{\frac{1}{k}}. \]
    Taking the limit of $K$ to $\infty$ we obtain the result.
  \end{proof}
\end{mytheo}

\begin{myexem}
  \label{exem:simple3}
  Suppose we apply Algorithm~\ref{algo:prod} to Example~\ref{exem:simple2}
  and the algorithm chooses $p_0(x) = \sum_\alpha c_\alpha x^\alpha$.
  The start of the sequence produced depends on the order between
  the coefficients $c_{(2d,0,0)},\allowbreak c_{(0,2d,0)},\allowbreak c_{(0,0,2d)}$.
  If $c_{(2d,0,0)}$ is the largest then
  the $\G$-admissible left-infinite sequence found is
  \[ \ldots,1,2,3,1,2,3,1,2,3. \]

  The product $A_{\sigma_1} A_{\sigma_2} A_{\sigma_3} \cdots = A_3A_2A_1A_3A_2A_1 \cdots $
  is periodic and has an asymptotic growth rate $\rho(A_{\sigma_1}A_{\sigma_2}A_{\sigma_3})^{1/3}\allowbreak = 1$.
  Hence $1 \leq \cjsr$.
  We saw in Example~\ref{exem:simple2} that $\jsrsos = 1$ for any $d$.
  Therefore $\cjsr = 1$.
\end{myexem}


\subsection{Refinement of the algorithm}
As a consequence of Theorem~\ref{theo:third} and Definition~\ref{def:cjsr},
\[ \jsrsos \leq [\mdin]^{\frac{1}{2d}} \cjsr. \]
However it is easy to see that $\mdin \geq \rho(A(\G))$ so
Theorem~\ref{theo:second} is a stronger statement.  This is somewhat
surprising, since Algorithm~\ref{algo:prod} uses more information than
the power iteration of \Sect\eqref{sec:quality} (i.e., the dual
solution of the SOS program), so we should be able to improve on it.


Indeed, we give now a generalization of Algorithm~\ref{algo:prod}.
The generalized algorithm has an integer parameter $l$.  When $l = 1$,
it reduces to Algorithm~\ref{algo:prod} which has a guarantee of
accuracy with $\mdin$ given by Theorem~\ref{theo:third}.  Increasing
$l$ increases the computation time of the generalized algorithm to
produce a sequence of the same length but it also improves the
guarantee.  With $l \to \infty$, it has the same guarantee except that
$\rho(A(G))$ replaces $\mdin$; see Theorem~\ref{theo:fourth}.
This section therefore provides an alternative proof of
Theorem~\ref{theo:second}.

Consider the example illustrated by Figure~\ref{fig:paths}.
Algorithm~\ref{algo:prod} considers all the \arcs{} ending at $v_1$ and chooses one of them.
This is what is illustrated by Figure~\ref{fig:path1}.
Then it denotes the start of the chosen \arc{} as $v_2$ and considers all the \arcs{} ending at $v_2$.
Instead of doing so, the algorithm could postpone the choice of $v_2$ and look at all the pairs $(v_3,v_2,\sigma_2),(v_2,v_1,\sigma_1)$ or equivalently, all the paths of length 2 whose last node is $v_1$ and choose one of them based on the dual variables of the edges $(v_3,v_2,\sigma_2)$ ignoring the dual variables of the edges $(v_2,v_1,\sigma_1)$.
The intuition behind this is that even if for large enough $d$ the accuracy is arbitrarily good,
for small $d$, we should not follow the dual variables ``blindly''.
Building the sequence with a ``larger horizon'' is a way to be less affected by the possible imprecision of the dual variables.


\begin{figure}[!ht]
  \begin{subfigure}{0.48\textwidth}
  \centering
    \begin{tikzpicture}
      \SetGraphUnit{1.2}
      \GraphInit[vstyle=Dijkstra]
      \SetVertexNoLabel
      \SetUpEdge[style={->}]
      \Vertex[NoLabel=false,L=$v_1$]{0}
      \NOWE(0){4};
      \NOEA(0){5};
      \SO[NoLabel=false,L=$v_0$](0){11};
      \tikzset{VertexStyle/.append style = {dashed}}
      \SOWE(0){10};
      \SOEA(0){12};
      \Edge[label=?](4)(0);
      \Edge[label=?](5)(0);
      \Edge(0)(11);
      \tikzset{EdgeStyle/.append style = {dashed}}
      \Edge(0)(10);
      \Edge(0)(12);
    \end{tikzpicture}
    \caption{One iteration of Algorithm~\ref{algo:prod}.}
    \label{fig:path1}
  \end{subfigure}
  \begin{subfigure}{0.48\textwidth}
  \centering
    \begin{tikzpicture}
      \SetGraphUnit{1.2}
      \GraphInit[vstyle=Dijkstra]
      \SetVertexNoLabel
      \SetUpEdge[style={->}]
      \Vertex[NoLabel=false,L=$v_1$]{0}
      \NOWE(0){4};
      \NOEA(0){5};
      \NOWE(4){6};
      \NO  (4){7};
      \NOEA(4){8};
      \NOEA(5){9};
      \SO[NoLabel=false,L=$v_0$](0){11};
      \tikzset{VertexStyle/.append style = {dashed}}
      \SOWE(4){1};
      \EA(0){2};
      \EA(2){3};
      \SOWE(0){10};
      \SOEA(0){12};
      \Edge[label=?](6)(4);
      \Edge[label=?](7)(4);
      \Edge[label=?](8)(4);
      \Edge[label=?](8)(5);
      \Edge[label=?](9)(5);
      \Edge[label=?](4)(0);
      \Edge[label=?](5)(0);
      \Edge(0)(11);
      \tikzset{EdgeStyle/.append style = {dashed}}
      \Edge(4)(1);
      \Edge(5)(2);
      \Edge(5)(3);
      \Edge(0)(10);
      \Edge(0)(12);
    \end{tikzpicture}
    \caption{One iteration of Algorithm~\ref{algo:prodl} with $l = 2$.}
    \label{fig:path2}
  \end{subfigure}
  \caption{Comparison between Algorithm~\ref{algo:prod} and Algorithm~\ref{algo:prodl}}
  \label{fig:paths}
\end{figure}

This idea can be generalized for any length $l$.
To shorten the notation we denote the $i$th node of a path $s$ as $s(i)$
and the $i$th label as $s[i]$.
\begin{algorithm}[!ht]
  \caption{Generates sequences of large asymptotic growth using paths of length $l$.}
  \label{algo:prodl}
  \begin{algorithmic}
    \STATE{Pick an arbitrary polynomial $p_0(x) \in \inte(\Sh)$}
    \STATE{Pick an \arc{} $(v_1,v_0,\sigma_0) \in \Arcs$ such that $\pE_{v_1v_0\sigma_0}[p_0(x)] > 0$}
    \STATE{Set $p_1(x) \leftarrow p_0(x)$}
    \FOR{$k = 1,2,\ldots$}
    \STATE $(v_{k+l},\sigma_{k+l-1},\ldots,\sigma_k,v_k) \leftarrow$
    \STATE $\hfill \argmax_{s \in \Arcs_l,s(l+1)=v_k} \pE_{s(1)s(2)s[1]}[p_k(A_s x)]$
    \STATE $p_{k+1} \leftarrow p_k(A_{\sigma_k} \cdots A_{\sigma_{k+l-1}} x)$
    \ENDFOR
  \end{algorithmic}
\end{algorithm}

Lemma~\ref{lemm:maxdegpath} is generalized into the following Lemma.

\begin{mylemm}
  \label{lemm:maxdegpathl}
  Consider \defAGe{}.
  For any positive integer $d$ and $l$,
  using Program~\ref{prog:dual} with any $\gamma < \jsrsos$,
  Algorithm~\ref{algo:prodl} with paths of length $l$ produces a $\G^\Tr $-admissible sequence
  $(v_1,v_0,\sigma_0), (v_2,v_1,\sigma_1), \ldots$ %
  that satisfies the following inequality for all $k \geq 1$ that is a multiple of $l$:
  \begin{equation*}
    \theta_k
    \geq \frac{\gamma^{2dl}}{\din_l(v_{k-l+1})}
    \theta_{k-l}.
  \end{equation*}
  where for a node $v \in \Nodes$, $\din_l(v)$ is the number of paths of length $l$ ending at $v$.
  \begin{proof}
    For any integer $k \geq 1$ multiple of $l$, since $p_0(x)$ is SOS, $p_{k-l+1}(x),\ldots,p_{k}(x)$ are all SOS.
    Therefore, by \eqref{eq:dualconsexp},

    \begin{multline*}
      \sum_{\substack{s \in \Arcs_l,\\s(l+1) = v_{k-l+1}}} \pE_{s(1)s(2)s[1]}[p_{k-l+1}(A_{s} x)]\displaybreak[3]\\
      \geq \gamma^{2d} \sum_{\substack{s \in \Arcs_{l-1},\\s(l) = v_{k-l+1}}} \pE_{s(1)s(2)s[1]}[p_{k-l+1}(A_{s} x)]\displaybreak[3]\\
      \vdots\qquad\qquad\qquad\qquad\qquad\qquad\displaybreak[3]\\
      \geq \gamma^{2dl} \theta_{k-l}\hspace{5.21cm}
    \end{multline*}
    Since the first expression has $\din_l(v_{k-l+1})$ positive terms and
    Algorithm~\ref{algo:prodl} picks the term with highest value, the result follows.
  \end{proof}
\end{mylemm}


Theorem~\ref{theo:third} is generalized in the following Theorem.

\begin{mytheo}
  \label{theo:fourth}
  Consider \defAGie.
  For any positive integers $d$ and $l$
  and any $\epsilon > 0$,
  using Program~\ref{prog:dual} with any $\gamma < \jsrsos$,
  Algorithm~\ref{algo:prodl} with paths of length $l$ produces a $\G^\Tr $-admissible sequence
  $(v_1,v_0,\sigma_0),\allowbreak(v_2,v_1,\sigma_1), \ldots$
  such that
  \begin{equation}
    \lim_{k \to \infty} \|A_{\sigma_1} \cdots A_{\sigma_k}\|_2^{\frac{1}{k}} \geq \frac{\gamma}{\|[A(\G)^\Tr ]^l\|_\infty^{\frac{1}{2dl}}}.
  \end{equation}
  \begin{proof}
    Consider the sequence produced by Algorithm~\ref{algo:prodl}.
    Since the set of \arcs{} $\Arcs$ is finite,
    there must be an \arc{} $(\bar{u},\bar{v},\bar{\sigma}) \in \Arcs$ that appears infinitely many times in the sequence
    at multiples of $l$.
    Let $k_1$ be the smallest multiple of $l$ such that $(v_{k_1+1},v_{k_1},\sigma_{k_1}) = (\bar{u},\bar{v},\bar{\sigma})$ and
    let $q(x) = p_{k_1+1}(x)$.
    Since the matrices of $\A$ are invertible and $p_0(x) \in \inte(\Sh)$, we know that $q(x) \in \inte(\Sh)$.

    For an arbitrarily large integer $K$, there exists an multiple of $l$ $k \geq K$ such that $(v_{k_1+k+1},v_{k_1+k}, \sigma_{k_1+k}) = (\bar{u},\bar{v},\bar{\sigma})$.
    Let $s_K = (\sigma_{k_1+k}, \ldots, \sigma_{k_1+1})$.
    By the definition of the matrix norm $\|\cdot\|_\infty$, $\max_{v \in \Nodes} \din_l(v) = \|[A(\G)^\Tr ]^l\|_\infty$.
    Theorefore, by Lemma~\ref{lemm:maxdegpathl}, we have
    \begin{equation}
      \label{eq:cyclingl}
      \pE_{\bar{u}\bar{v}\bar{\sigma}}[q(A_{s_K} x)] \geq \frac{\gamma^{2dk}}{\|[A(\G)^\Tr ]^{l}\|_\infty^{k/l}} \pE_{\bar{u}\bar{v}\bar{\sigma}}[q(x)].
    \end{equation}
    By Lemma~\ref{lemm:sosbound}, there exists a constant $\beta$ that does not depend on $A_{s_K}$ such that%
    \[ \beta \|A_{s_K}\|_2^{2d}q(x) - q(A_{s_K} x) \text{ is SOS}. \]
    Therefore by \eqref{eq:cyclingl},
    \[ \beta \|A_{s_K}\|_2^{2d}\pE_{\bar{u}\bar{v}\bar{\sigma}}[q(x)] \geq \frac{\gamma^{2dk}}{\|[A(\G)^\Tr ]^{l}\|_\infty^{k/l}} \pE_{\bar{u}\bar{v}\bar{\sigma}}[q(x)]. \]
    Hence,
    \[ \frac{\gamma}{[\beta^{\frac{1}{k}} \|[A(\G)^\Tr ]^{l}\|_\infty^{1/l}]^{\frac{1}{2d}}} \leq \|A_{s_K}\|_2^{\frac{1}{k}}. \]
    Taking the limit of $K$ to $\infty$ we obtain the result.
  \end{proof}
\end{mytheo}

Since $\lim_{l \to \infty} \|[A(\G)^\Tr ]^l\|_\infty^{\frac{1}{l}} = \rho(A(\G)^\Tr ) = \rho(A(\G))$, Theorem~\ref{theo:second} is a consequence of Theorem~\ref{theo:fourth}.


\subsection{Producing lower bounds}
\label{sec:lb}
By definition of the CJSR, the asymptotic growth rate of the norm of the product of any $\G$-admissible (or $\G^\Tr $-admissible) sequence of matrices gives a lower bound for the CJSR.
In particular the sequence produced by Algorithm~\ref{algo:prodl} provides a lower bound for the CJSR.

If there are two integers $\bar{k},k$ such that the sequence after $\bar{k}$ is periodic of period $k$,
the asymptotic growth rate of the norm is equal to the $k$th root of the spectral radius of the product of the matrices of one period.
This is due to the Gelfand's formula $\rho(A) = \lim_{k \to \infty} \|A^k\|^{1/k}$.
From the same identity,
we see that the spectral radius of the product of the matrices of one $\G$-admissible cycle gives a lower bound for the CJSR.


More formally, for any natural $k$, we have
\[ \rho_k(\G, \A) \leq \cjsr \]
where
\[ \rho_k(\G, \A) = \max \big\{\, \rho(A_c)^{1/k} : c \in \G_k, c \text{ is a cycle}\,\big\}. \]
This is the generalization to the constrained case of the first inequality of the three members inequalities; see \cite[\Sect1.2.2.6]{jungers2009joint}.

To find lower bounds for the CJSR,
one could generate all the cycles of length smaller than some maximum length
and compute the spectral radius for all of them.
This brute force approach is not scalable because the number of paths considered grows exponentially with the maximum length.
The exponential growth of the brute force approach is the reason why one should choose small $l$ for Algorithm~\ref{algo:prodl}.

Gripenberg~\cite{gripenberg1996computing} proposes a branch-and-bound algorithm that prunes the search using an a priori fixed absolute error.
Two other branch-and-bound variant exists:
the balanced complex polytope algorithm \cite{guglielmi2008algorithm}
and the invariant conitope algorithm~\cite{jungers2014lifted}.

These algorithms can also be used to produce a $\G$-admissible sequence of matrices of high asymptotic growth rate
by reproducing the cycles of high spectral radius infinitely.
The advantage of Algorithm~\ref{algo:prodl} is that it provides a guarantee of accuracy given in Theorem~\ref{theo:fourth}.
Algorithm~\ref{algo:prodl} provides at the same time unstable trajectories and lower bounds of guaranteed accuracy.

We can compute lower bounds using the upper bound provided by Program~\ref{prog:primal} and Corollary~\ref{cor:mix}
but in practice the trajectories are periodic after some time $\bar{k}$ so we are able to compute
much better lower bounds than the pessimistic bound provided by Corollary~\ref{cor:mix}.
This is shown by the following example.

\begin{myexem}
  \label{exem:run4}
  We first analyse the behaviour of Algorithm~\ref{algo:prod}
  on our running example; see Example~\ref{exem:run1}, Example~\ref{exem:run2} and Example~\ref{exem:run3}.
  We use the recent CSSystem toolbox \cite{cambier2015cssystem}.
  For $2d = 2,4,6,8$, the sequence can only contain the edges $(1,3,1)$, $(3,1,2)$ and $(3,3,1)$
  since the other edges have a dual variables of value 0.

  To find a cycle with good spectral radius, we generate the sequence for a fixed length and then look through all the cycles of the sequence.
  The $\G$-admissible cycle found depends on $p_0(x)$, but most of the time, it is
  \[ (3,1,2),(1,3,1),(3,1,2),(1,3,1),(3,3,1)^4 \]
  where the $k$ in exponent means that the edges is taken $k$ times.
  The 8th root of the corresponding spectral radius is 0.97289.

  Using Algorithm~\ref{algo:prodl} with $l = 3$, the produced sequence may include two consecutive edges that have a dual variables of value 0.
  Even for $2d = 2$, the Algorithm~\ref{algo:prodl} with $l = 3$ finds the $\G$-admissible cycle
  \[ (3,1,2),(1,3,1),(3,1,2),(1,2,3),(2,3,1),(3,3,1)^3 \]
  for which the 8th root of the corresponding spectral radius is 0.97482.
  We can show that it is equal to the CJSR, e.g. using \cite{guglielmi2008algorithm}.
  That is, it is a \emph{spectrum maximizing product} (s.m.p.).

  With $2d = 10$ and $2d = 12$, we have seen in Example~\ref{exem:run3} that $\pE_{123}$ and $\pE_{231}$ are not zero.
  For $2d = 12$, the s.m.p. is found for any $l$, initial node and $p_0(x)$.
  For $2d = 10$, despite the fact that $\pE_{123}$ and $\pE_{231}$ are not zero,
  the s.m.p. seems to be never found for $l = 1$.

  For $2d = 10$ and $l = 2$, it is found from the node 1 and 4\footnote{In practice, the dual variables of the edges $(3,4,4)$ and $(4,3,1)$ are close to zero but not zero. It is possible to start in node 4 even if the sequence will never go back in node 4.}
  with high probability from a random $p_0(x)$ or from $p_0(x)$ equal to the primal solution at this node.
  From node 3, it is not found using the primal solution of node 3 but it is sometimes found using a random $p_0(x)$.
  This shows that choosing the primal solution as $p_0(x)$ may not be the best choice.
  From node 2, it does not seem to be found for any $p_0(x)$.

  Instead of increasing $l$ and $d$, path-complete techniques
  \cite{ahmadi2014joint} can be used to improve the lower bounds
  provided by Algorithm~\ref{algo:prodl} and the upper bounds provided
  by Program~\ref{prog:primal}.  The idea is to replace the automaton
  $\G$ and the matrices of $\A$ by another automaton $\G'$ and other
  matrices $\A'$ but without changing the set of admissible products
  of matrices.  Typically, the new set $\A'$ contains products of the
  matrices of the initial set $\A$ and the automaton $\G'$ is designed
  so as to maintain the same set of admissible products.

  We tried two examples: The $T$-product lift
  \cite{philippe2015stability} and $M$-path-dependent lift
  \cite{lee2006uniform}.  Using the 2-product lift, the s.m.p. is
  found using $2d = 6$ and $l = 1$ and using the 2-path-dependent
  lift, the s.m.p. is found using $2d = 4$ and $l = 1$.
  This is concordant with the empirical observation reported in
  \cite{philippe2015stability}.
\end{myexem}


\section{Conclusions}

We have introduced a new technique, based on the dual solution of the
SOS approximation, to generate a sequence of matrices with asymptotic
growth rate close to the CJSR. In Theorem~\ref{theo:fourth}, we gave
an estimate of its performance. However, in practical examples, the
sequence produced achieves a growth rate which is often better than
this estimate.  This gives rise to the following question: Is it
possible to improve the guarantee both for the upper bounds and the
lower bounds?  To the best of our knowledge, there are currently two
bounds, Theorem~\ref{theo:matt} only depending on the dimension of the
matrices and Theorem~\ref{theo:second} only depending on the automaton.
Is there a bound that combines both in a more sophisticated way than
Corollary~\ref{cor:mix}?

Our algorithm seems to always produce periodic sequences so in
practice we can compute an asymptotic growth rate bound better than
what is provided by Corollary~\ref{cor:mix}.  This motivates other
questions with important algorithmic consequences: Does
Algorithm~\ref{algo:prodl} always produce a sequence that is periodic
after some $\bar{k}$?  Can we give a bound for $\bar{k}$ and the size
of the period depending on the dimension of the matrices and on the
automaton?

More generally, the techniques developed in this work, based on
generating ``bad'' trajectories for a dynamical system via the dual
solution of the natural convex problems used for analysis, naturally
extend to many other problems in systems theory. We are currently
exploring such possibilities.


\bibliographystyle{abbrv}

\bibliography{biblio}

\end{document}
